\subsection{Rectangular hooks in ordinary group cohomology}
This is an application of the known rectangular hook formula and
a known split-group motive decomposition, rather than a new
Kronecker rule.

\begin{finding}[\fid{F33}: a group-scheme realization with mixed degrees]\label{f:hookgeometry}
Let $n,d\ge1$, $N=nd$, $m=\min(n,d)$ and $0\le k<N$.
For $G=SL_{m,\F_q}$ and $\ell\nmid q$,
\[
 g((N-k,1^k),(d^n),(d^n))
   =\dim_{\Q_\ell}H^k(G_{\bar{\F}_q},\Q_\ell).
\]
More precisely there is an actual Weil decomposition
\begin{equation}\label{eq:hookweil}
 H^k(G_{\bar{\F}_q},\Q_\ell)=
 \bigoplus_{\substack{I\subseteq\{2,\ldots,m\}\\
                     \sum_{i\in I}(2i-1)=k}}
       \Q_\ell\left(-\sum_{i\in I}i\right).
\end{equation}
The degree can be mixed: for $m=8$,
\[
 H^{15}(SL_{8,\bar{\F}_q},\Q_\ell)
       =\Q_\ell(-8)\oplus\Q_\ell(-9).
\]
This realizes the numerical coefficients in ordinary cohomology;
it does not identify a nonstandard action or canonical basis.
\end{finding}

\begin{proof}
The known formula in IP, Proposition~6.4, is
\[
 \sum_{k=0}^{N-1}g((N-k,1^k),(d^n),(d^n))t^k
       =\prod_{i=2}^m(1+t^{2i-1}).
\]
For the geometric input, Merkurjev gives the canonical split
special-linear-group motive decomposition over a perfect field
\cite[\S5, Corollary~5.5, p.~8; notation on p.~7]{MerkurjevSL}.
In his notation a sequence
$1\le j_1<\cdots<j_a\le m-1$ contributes
$\mathbf1(d+a)[2d+a]$, with $d=\sum j_r$.
Reindexing by $I=\{j_1+1,\ldots,j_a+1\}$ and
$s_I=\sum_{i\in I}i$ gives
\[
 M(SL_m)_{\Q}
       \simeq\bigoplus_{I\subseteq\{2,\ldots,m\}}
                 \mathbf1(s_I)[2s_I-|I|].
\]
The displayed known motive theorem is invoked, not reproved.
Finite fields satisfy its perfect-field hypothesis.
Under ordinary cohomological $\ell$-adic realization, the summand
$\mathbf1(s)[a]$ contributes $\Q_\ell(-s)$ in degree $a$.
The decomposition is defined over $\F_q$, so Frobenius on that
summand is exactly $q^s$. This proves \eqref{eq:hookweil}.
Taking dimensions yields the product above, hence the required
coefficient identity.

For a summand in degree $k=2s_I-|I|$, its weight is
$2s_I=k+|I|$. Different subset cardinalities can therefore give
different weights in the same degree. At $m=8,k=15$, the only
subsets are $\{8\}$ and $\{2,3,4\}$. A two-element subset has
even degree, while four or more elements have minimum degree
$3+5+7+9=24$. Among three-element subsets the minimum 15 is
attained only by $\{2,3,4\}$. Their sums are eight and nine,
giving the stated Tate summands, weights 16 and 18, and traces
$q^{8r}+q^{9r}$ over $\F_{q^r}$.
\end{proof}

The corresponding hook coefficient is two for every rectangle
whose minimum side is eight. A single Tate twist cannot make the
two displayed weights equal. This is consistent with \fid{F20},
whose pure Satake multiplicity motive is a different realization
of the same integer. The affine group $G$ supplies ordinary
cohomology; no assertion of projectivity or purity of each
cohomological degree is needed.

% Separate paragraph for Part III.
\paragraph{\fid{F33} in Grothendieck's language.}
The geometric object is the ordinary split group scheme $SL_m$.
Its rational Tate-motive decomposition supplies actual Weil
summands and scheme Frobenius, so the equality of dimensions
has a precise realization. The relevant frameworks are SGA~3's
split group schemes and the explicitly cited modern motive
theorem. The generator degree $2i-1$ and weight $2i$ remain
different gradings. No manuscript leaf is claimed to contain
this hook formula or a nonstandard basis comparison.

% New bibliography key:
% MerkurjevSL: A. S. Merkurjev, Motivic decomposition of certain
% special linear groups, inspected author PDF, Section5,
% Corollaries5.3 and5.5, pp7--8.
% https://www.math.uni-bielefeld.de/LAG/man/530.pdf
